\section*{Technical supplement}
\phantomsection
\label{supp:start}
\addcontentsline{toc}{section}{Technical supplement}

The main article is a synthesis of the current evidence. This supplement is
the detailed record needed to examine how that evidence was obtained.
It retains negative experiments, numerical settings, intermediate hypotheses,
and corrections that would interrupt the main narrative.

\begin{description}
\item[Appendix~\ref{supp:methods}, page~\pageref{supp:methods}.]
Numerical methods: orbit identity, shooting, Floquet calculations, saddle
reconstruction, and acceptance tests.
\item[Appendix~\ref{supp:experiments}, page~\pageref{supp:experiments}.]
The experimental sequence and the remaining 24 scientific figures.
References to the ten main figures retain their original labels.
\item[Appendix~\ref{supp:symbolic}, page~\pageref{supp:symbolic}.]
Source transcription, operational symbols, unsuccessful center searches,
and the remaining reinjection tests.
\item[Appendix~\ref{supp:interpretation}, page~\pageref{supp:interpretation}.]
How interpretations changed as controls and independent tests accumulated.
\item[Appendix~\ref{supp:reviewers}, page~\pageref{supp:reviewers}.]
The original reviewer concerns and the evidence still required to address them.
\end{description}

Numbered experiments refer to the frozen manifests and machine-readable
receipts in the repository. Chronological statements describe the state of
the project at that experiment, not necessarily its final conclusion.
The main article, especially Table~\ref{tab:current-claims}, takes precedence
over an interpretation explicitly superseded later in this record.
Supplementary figures use the prefix S; asset filenames keep their stable
historical identifiers so their generation records remain traceable.

\clearpage
